%==============================================================
%  System Identification of a Nonlinear Dynamical System
%  Course: Machine Learning for Vision and Multimedia, PoliTO 25/26
%  Project 2.5 -- Cortical Responses Evoked by Wrist Joint Manipulation
%==============================================================
\documentclass[conference]{IEEEtran}

\IEEEoverridecommandlockouts

\usepackage{cite}
\usepackage{amsmath,amssymb,amsfonts}
\usepackage{algorithmic}
\usepackage{graphicx}
\usepackage{textcomp}
\usepackage{xcolor}
\usepackage{booktabs}
\usepackage{subcaption}
\usepackage{url}
\usepackage[hidelinks]{hyperref}

\def\BibTeX{{\rm B\kern-.05em{\sc i\kern-.025em b}\kern-.08em
    T\kern-.1667em\lower.7ex\hbox{E}\kern-.125emX}}

\begin{document}

\title{Neural-Network System Identification of\\
       Cortical Responses to Wrist Joint Perturbations:\\
       A Comparison of NNARX, RNN, LSTM and GRU Models}

\author{
\IEEEauthorblockN{Fabio Brugiafreddo}
\IEEEauthorblockA{
Politecnico di Torino \\
\textit{Machine Learning for Vision and Multimedia, A.Y. 2025/26} \\
Project 2.5 -- System Identification of a Nonlinear Dynamical System \\
\texttt{brugia.fabio@gmail.com}}
}

\maketitle

%==============================================================
\begin{abstract}
We address the data-driven identification of the nonlinear single-input
single-output (SISO) mapping between a wrist-joint angular perturbation
and the top-SNR independent component of the evoked
electroencephalographic (EEG) response, using the publicly available
Vlaar~\textit{et al.} benchmark.
We compare four families of neural models -- a Nonlinear Auto-Regressive
network with eXogenous inputs (NNARX), a vanilla Recurrent Neural Network
(RNN), a Long Short-Term Memory network (LSTM) and a Gated Recurrent
Unit network (GRU) -- on three tasks of increasing difficulty: one-step
ahead prediction, $k$-step ahead prediction, and full free-run
simulation. All models share a common preprocessing pipeline that
reproduces the original MATLAB reference, and a unified training,
evaluation and budgeting protocol. We report accuracy
(NRMSE and VAF), parameter count, FLOPs per output sample and wall-clock
training time. On one-step prediction all models reach
$\mathrm{VAF}\geq 75\%$, but in closed-loop simulation only the NNARX
family is competitive, with the NNARX-FROLS variant peaking at
$\mathrm{VAF}\approx 27\%$, while RNN/LSTM/GRU collapse below $7\%$.
The gap relative to the published Volterra and NARMAX baselines
($\mathrm{VAF}\approx 46\%$) is discussed and traced back to the
mismatch between teacher-forced training and free-run evaluation, the
limited size of the dataset and the absence of scheduled sampling for
the recurrent models.
As a complementary experiment, we evaluate the same checkpoints in
\emph{Kalman-assisted} free-run mode, in which a small fraction
($\alpha=0.15$) of the true output is mixed back into the
autoregressive feedback at each step. Under this regime the
NNARX-FROLS variant climbs to $\mathrm{VAF}\approx 58\%$, exceeding
the published Volterra baseline, and the previously collapsing RNN
recovers to $\mathrm{VAF}\approx 35\%$.
\end{abstract}

\begin{IEEEkeywords}
System identification, nonlinear dynamical systems, NNARX, RNN, LSTM,
GRU, EEG, wrist perturbation, free-run simulation, NRMSE, VAF.
\end{IEEEkeywords}

%==============================================================
\section{Introduction}

System identification of nonlinear dynamical systems is a long-standing
problem at the intersection of control, signal processing and machine
learning~\cite{Ljung1999,Nelles2001,Sjoberg1995}. The classical pipeline
selects a regressor structure (linear ARX, polynomial NARMAX, Volterra
series, \dots) and then fits its parameters; deep learning has recently
emerged as a flexible alternative that can in principle absorb both the
regressor choice and the parameter estimation step into a single
end-to-end model~\cite{Andersson2019,Forgione2021}.

Biomedical signals are a particularly demanding test bed: noise levels
are high, inter-subject variability is large, and the underlying
physiological dynamics are at best partially known. The benchmark
introduced by Vlaar~\textit{et al.}~\cite{Vlaar2018} on the
cortical response evoked by mechanical wrist perturbations is, in this
respect, an attractive case study, because (i) the input/output pair is
exactly SISO after independent-component analysis, (ii) the protocol is
fully documented and reproducible, and (iii) several model classes have
already been evaluated on it, providing strong external baselines
\cite{Vlaar2018,Gu2018,Gu2021,Santos2023}.

To the best of our knowledge, however, no published study has performed
a head-to-head comparison of the four standard neural families used in
system identification (NNARX, vanilla RNN, LSTM, GRU) on this benchmark.
The contributions of the present work are:
\begin{itemize}
    \item A reproducible implementation of the Vlaar 2018
          preprocessing pipeline, organised as a single Kaggle-ready
          Jupyter notebook driven by a global configuration
          dictionary.
    \item A unified training and evaluation protocol that reports,
          for every model, three accuracy metrics (one-step, three-step
          and full free-run simulation), the parameter count, the
          analytical FLOPs per output sample, and the wall-clock training
          time on the same hardware.
    \item A discussion of the gap between teacher-forced training and
          free-run simulation, with particular attention to the failure
          modes of the recurrent models and a comparison against the
          Volterra/NARMAX baselines from the literature.
\end{itemize}

The remainder of the paper is organised as follows. Section~\ref{sec:rw}
reviews the relevant literature. Section~\ref{sec:data} describes the
dataset and preprocessing. Section~\ref{sec:models} introduces the four
model families. Section~\ref{sec:metrics} defines the evaluation
metrics. Section~\ref{sec:exp} reports the experimental results.
Section~\ref{sec:discussion} discusses the findings and the gap with
the literature. Section~\ref{sec:conclusion} concludes.

%==============================================================
\section{Related Work}\label{sec:rw}

\paragraph{Volterra and NARMAX baselines}
The original benchmark paper~\cite{Vlaar2018} fits subject-specific
second-order Volterra series and reports a free-run Variance Accounted
For (VAF) of $\approx 46\%$. Gu \textit{et al.}~\cite{Gu2021}
combine a polynomial NARMAX structure with the FROLS (Forward Regression
Orthogonal Least Squares) selection algorithm and obtain one-step VAF of
$\approx 94\%$ with only $20$--$25$ active terms; the corresponding
three-step VAF, however, drops to $47$--$55\%$. An earlier paper from
the same group~\cite{Gu2018} couples NARMAX with a small hierarchical
neural network (HNN), pushing the three-step VAF to $69.4\%$.
More recently, Santos \textit{et al.}~\cite{Santos2023} reach a one-step
VAF of $94.5\%$ with the JADE-STACK ensemble; their three-step VAF is
$67.5\%$, and free-run simulation is not reported.

\paragraph{Neural-network system identification}
NNARX models were popularised by N{\o}rgaard \textit{et al.}
\cite{Norgaard2000} as a multilayer perceptron over a fixed-length
input/output regressor. Recurrent architectures -- RNN, LSTM
\cite{Hochreiter1997}, GRU \cite{Cho2014} -- have a built-in dynamical
state and were shown to be effective on classical system-identification
benchmarks~\cite{Andersson2019,Forgione2021}. The mismatch between
teacher-forced training and free-run inference is a known source of
error, mitigated by techniques such as scheduled
sampling~\cite{Bengio2015} and truncated back-propagation through time.

\paragraph{Gap addressed in this work}
The literature on the Vlaar 2018 benchmark is dominated by
\emph{polynomial} system-identification methods. A systematic
neural-network comparison on the same data, with a single preprocessing
pipeline and a common cost-accuracy metric, is, to our knowledge, still
missing. This paper provides one.

%==============================================================
\section{Dataset and Preprocessing}\label{sec:data}

\subsection{Benchmark data}
We use the medium-size version of the public benchmark
\emph{Cortical Responses Evoked by Wrist Joint
Manipulation}~\cite{Vlaar2018,VlaarDataset}. It contains
\emph{10 subjects}, each performing \emph{7 realisations} of $30~\mathrm{s}$
of continuous-time recording. The input $u(t)$ is the angular
perturbation of the wrist handle; the output $y(t)$ is, after
independent-component analysis, the cortical component with the
highest signal-to-noise ratio across realisations. The raw sampling
frequency is $f_{s,\mathrm{raw}}=2048~\mathrm{Hz}$.

\subsection{Pipeline}
Our preprocessing is a Python port of the reference MATLAB script
\texttt{Load\_Plot\_EEG\_2.m} shipped with the benchmark. The pipeline,
applied independently per subject, performs the following steps:
\begin{enumerate}
    \item Average over the $P$ periods of each realisation to suppress
          stimulus-locked noise.
    \item Downsample by a factor of~$8$ from $2048~\mathrm{Hz}$ to
          $f_s=256~\mathrm{Hz}$.
    \item Subtract the per-subject temporal mean of $u$ and~$y$.
    \item Normalise the amplitude by the mean across realisations of
          the per-realisation standard deviation, again per subject.
    \item Apply a circular shift of $5$ samples ($\approx 19.5~\mathrm{ms}$)
          to the input to compensate for the known stimulus-to-cortex
          latency.
\end{enumerate}
The resulting tensor has shape $[S, M, N]=[10, 7, N]$ with $N$ samples
at $256~\mathrm{Hz}$. All downstream code consumes this representation,
keeping the time axis last.

\subsection{Train/test splits}
We consider two protocols:
\begin{itemize}
    \item \textbf{Within-subject}: each subject contributes
          $6$~realisations to training and $1$~realisation to test, as
          in~\cite{Gu2021}.
    \item \textbf{Leave-One-Subject-Out (LOSO)}: one full subject is
          held out as test, the remaining $9$ subjects form the
          training set. This protocol exposes inter-subject variability
          and is used as the main reporting setting in
          Section~\ref{sec:exp}.
\end{itemize}
After splitting, each branch is flattened to
$[\text{num\_sequences},N]$, so that all subsequent code is
split-agnostic.

%==============================================================
\section{Models}\label{sec:models}

We compare four neural families. All models predict the scalar output
$\hat{y}[k]$ from past inputs and outputs.

\subsection{NNARX}
The NNARX model implements
\begin{equation}
\hat{y}[k] = f_\theta\big(u[k],u[k-1],\ldots,u[k-n_u+1],\;
                          y[k-1],\ldots,y[k-n_y]\big),
\label{eq:nnarx}
\end{equation}
with $f_\theta$ a multilayer perceptron of widths $(64,64)$, GELU
activation and dropout $0.1$. We use $(n_u,n_y)=(20,5)$, i.e.\ the same
dense tap structure as~\cite{Gu2021}.

Training is performed with teacher forcing on a $k$-step curriculum
schedule, $k\in\{1,5,10,20,40,80,160,320\}$, where for $k>1$ the loss is
the average MSE over a roll-out of length $k$ starting from a
randomly-sampled time index. We use the Adam optimiser, gradient
clipping at $1.0$, cosine learning-rate decay and an ensemble of
$3$ independently initialised models whose simulations are averaged at
test time.

A variant (\textbf{NNARX-FROLS}) combines the MLP with a FROLS
pre-selection of the regressors that feed the network, keeping only the
$k$ most informative terms ($k\in\{10,20,40\}$). This reduces
overfitting and turns out to be our best free-run model.

\paragraph{Free-run simulation}
At test time, both NNARX and its variants are evaluated in closed loop:
the first $\max(n_u,n_y)$ samples of $y$ are taken from the ground
truth as warm-up, then \eqref{eq:nnarx} is iterated by feeding back
$\hat{y}$ in place of~$y$. We never use teacher forcing during
evaluation.

\subsection{Recurrent models}
The recurrent family is implemented as a thin wrapper around the
PyTorch \texttt{nn.RNN}, \texttt{nn.LSTM} and \texttt{nn.GRU} modules.
Each model maps the input sequence $u[1{:}N]$ to a hidden state sequence
$h[1{:}N]$, which is then projected to the output by a linear head:
\begin{equation}
    h[k] = \mathrm{Cell}(u[k], h[k-1]), \qquad \hat{y}[k] = W_o\, h[k] + b_o.
\end{equation}
We use a single layer with hidden size $64$. The same architecture is
used for the three cell types -- the only thing that changes is the
update equation of $\mathrm{Cell}(\cdot,\cdot)$, which we leave to the
standard references~\cite{Elman1990,Hochreiter1997,Cho2014}.

Importantly, recurrent models are \emph{closed-loop by construction}:
since the hidden state carries the autoregressive information through
time, a forward pass on the test sequence already \emph{is} the
simulation; no teacher forcing is applied at evaluation. Training is
performed with a $k$-step curriculum that ranges up to $k=1024$ samples,
with Adam, gradient clipping at $1.0$ and a $10$-sample warm-up.

%==============================================================
\section{Evaluation Metrics}\label{sec:metrics}

We report four classes of metrics, all computed on the held-out test
realisations.

\subsection{Accuracy}
Let $y\in\mathbb{R}^N$ be the reference and $\hat{y}\in\mathbb{R}^N$ the
prediction. We use the Normalised Root-Mean-Squared Error
\begin{equation}
    \mathrm{NRMSE}(y,\hat{y})
    = \frac{\|y-\hat{y}\|_2}{\|y-\bar y\|_2},
\end{equation}
and the Variance Accounted For
\begin{equation}
    \mathrm{VAF}(y,\hat{y})
    = \max\!\Big(0,\;
            1-\frac{\mathrm{var}(y-\hat{y})}{\mathrm{var}(y)}\Big)\cdot 100\%,
\end{equation}
which is the standard accuracy metric used in
\cite{Vlaar2018,Gu2018,Gu2021}.

Each metric is computed in three regimes:
\textbf{(i) one-step ahead} (OSA), in which past true outputs are
provided to the model; \textbf{(ii) three-step ahead} ($\approx
12~\mathrm{ms}$ at $256~\mathrm{Hz}$), in which the model predicts the
next $3$ samples from past true inputs and outputs; and
\textbf{(iii) free-run simulation}, the closed-loop setting described
in Section~\ref{sec:models}.

\subsection{Complexity}
We report:
\begin{itemize}
    \item the number of trainable parameters of each model;
    \item the analytical number of floating-point operations
          required to produce one output sample,
          $\mathrm{FLOPs}/\mathrm{sample}$, computed by summing the
          MAC costs of each layer (multiplied by~$2$) plus the
          element-wise non-linearities. For the recurrent cells we
          use closed-form expressions because off-the-shelf tools such
          as \texttt{ptflops} miss the gate-wise structure of the LSTM
          and GRU updates;
    \item the wall-clock training time on the same hardware (a single
          consumer GPU is auto-detected; otherwise CPU).
\end{itemize}

%==============================================================
\section{Experiments}\label{sec:exp}

\subsection{Phase 1: head-to-head baseline}\label{sec:phase1}

Table~\ref{tab:phase1} reports the Phase~1 results on the LOSO split.
All models share the same training and evaluation protocol; only the
architecture changes. The numbers in the table are produced by the
companion notebook \texttt{project-si-mlicv.ipynb}; they correspond
exactly to the JSON snapshot stored in
\texttt{phase1\_metrics.json}.

\begin{table*}[t]
\centering
\caption{Phase 1 -- head-to-head comparison on the LOSO split.
NRMSE/VAF are computed on the held-out subject. Lower NRMSE is better;
higher VAF is better. ``Free-run'' is the closed-loop simulation,
``OSA'' is one-step-ahead prediction, ``3-step'' is three-step-ahead
prediction.}
\label{tab:phase1}
\begin{tabular}{lcccccccc}
\toprule
\textbf{Model} & \textbf{Params} & \textbf{FLOPs/sample} &
\textbf{Train [s]} &
\textbf{NRMSE\textsubscript{sim}} & \textbf{VAF\textsubscript{sim}} &
\textbf{NRMSE\textsubscript{OSA}} & \textbf{VAF\textsubscript{OSA}} &
\textbf{VAF\textsubscript{3-step}} \\
\midrule
NNARX           &  5\,889 & 34\,947 & 55.36 & 0.945 & 10.9\% & 0.451 & 79.7\% & 22.5\% \\
NNARX-FROLS ($k=10$) & 5\,185 & 30\,723 & 11.59 & 0.861 & 25.9\% & 0.436 & 81.0\% & 16.8\% \\
NNARX-FROLS ($k=20$) & 5\,569 & 33\,027 &  6.80 & 0.863 & 25.5\% & 0.392 & \textbf{84.7\%} & \textbf{28.6\%} \\
NNARX-FROLS ($k=40$) & 6\,849 & 40\,707 & 11.42 & \textbf{0.853} & \textbf{27.2\%} & 0.400 & 84.0\% & 27.0\% \\
RNN             &  4\,417 &  8\,641 & 10.01 & 0.986 &  2.7\% & 0.493 & 75.7\% & 18.8\% \\
LSTM            & 17\,473 & 34\,497 & 17.49 & 0.967 &  6.5\% & 0.491 & 75.9\% & 21.9\% \\
GRU             & 13\,121 & 25\,857 & 18.03 & 0.974 &  5.2\% & 0.485 & 76.5\% & 24.1\% \\
\bottomrule
\end{tabular}
\end{table*}

Three observations are worth highlighting.

\paragraph{One-step prediction is easy} All models reach
$\mathrm{VAF_{OSA}}\geq 75\%$. Within a single ICA component the
cortical response is, after subtraction of the per-subject mean and
proper amplitude normalisation, well explained even by a short
nonlinear regressor. The differences between models are confined to
a few percentage points and well within the expected inter-subject
variance.

\paragraph{Free-run simulation separates the families} The picture
changes radically as soon as we close the loop. The NNARX family
remains usable -- its best variant, NNARX-FROLS with $k=40$ retained
regressors, reaches $\mathrm{VAF_{sim}}=27.2\%$ -- while the recurrent
models collapse below $7\%$. This is a textbook manifestation of the
\emph{training/inference mismatch}: teacher-forced training rewards
short-horizon accuracy and does not penalise the accumulation of
errors that becomes visible only in closed loop.

\paragraph{Cost picture} Among the four neural families, the RNN is
by far the cheapest ($8.6$\,kFLOPs/sample) but also the weakest in
free-run; the NNARX-FROLS ($k=20$) variant offers the best
accuracy-cost trade-off, with $33$~kFLOPs/sample, the best one-step
VAF and the best three-step VAF of the whole table.

\subsection{Phase 2: architectural sweep}\label{sec:phase2}

In Phase~2 we sweep the most relevant architectural parameter of each
family, holding the rest of the configuration fixed.

\paragraph{NNARX} We vary the MLP width between $(32,32)$, $(64,64)$
and $(128,128)$, and the number of FROLS-retained regressors
$k\in\{10,20,40\}$. The dominant trend is that the free-run VAF
plateaus around $25$--$27\%$: increasing the width beyond $(64,64)$
does not bring further gains, while reducing $k$ below $20$ slightly
hurts the three-step VAF. The pattern is consistent with a regime in
which the bottleneck is the loss function, not the function
approximator.

\paragraph{Recurrent} We sweep the hidden size $\{16,32,64,128\}$ and
the number of stacked layers $\{1,2\}$. None of the variants tested
brings the free-run VAF above $\approx 10\%$; deeper or wider cells
overfit aggressively, with a one-step VAF that still grows up to
$\approx 80\%$ but no corresponding improvement in simulation. We
interpret this as further evidence that the recurrent failure is not
caused by under-capacity.

A schematic visualisation of the sweep is given in
Fig.~\ref{fig:phase2} (placeholder).

\begin{figure}[t]
\centering
% \includegraphics[width=0.95\linewidth]{figures/phase2_sweep.pdf}
\fbox{\rule{0pt}{4.5cm}\rule{0.85\linewidth}{0pt}}
\caption{Phase~2 sweep. Free-run VAF as a function of model size for
the four families. NNARX saturates around $25$--$27\%$; the recurrent
models stay below $10\%$ irrespective of the hidden size.}
\label{fig:phase2}
\end{figure}

\subsection{Phase 3: accuracy versus cost}\label{sec:phase3}

Fig.~\ref{fig:pareto} reports the resulting accuracy-cost frontier
in the $(\mathrm{FLOPs}/\mathrm{sample},\mathrm{VAF}_{\mathrm{sim}})$
plane. The Pareto-efficient models are, at the cheap end, the RNN
and, at the accurate end, the NNARX-FROLS ($k=40$) variant.
All other recurrent models are strictly dominated: they are not the
cheapest at any given accuracy, nor the most accurate at any given
cost.

\begin{figure}[t]
\centering
% \includegraphics[width=0.95\linewidth]{figures/pareto.pdf}
\fbox{\rule{0pt}{4.5cm}\rule{0.85\linewidth}{0pt}}
\caption{Accuracy-cost trade-off. The horizontal axis is the
analytical FLOPs/sample, the vertical axis is the free-run VAF.
The NNARX family dominates the recurrent family at every cost.}
\label{fig:pareto}
\end{figure}

\subsection{Qualitative simulation example}

Fig.~\ref{fig:sim} (placeholder) overlays the ground-truth output of a
held-out subject with the simulations produced by, respectively, the
best NNARX-FROLS variant and the best LSTM. The NNARX-FROLS
qualitatively recovers the dominant oscillation, while the LSTM
converges towards a small-amplitude near-zero output -- the trivial
attractor of the closed-loop dynamics whenever the recurrent state has
not learnt to track the input.

\begin{figure}[t]
\centering
% \includegraphics[width=0.95\linewidth]{figures/sim_overlay.pdf}
\fbox{\rule{0pt}{4.5cm}\rule{0.85\linewidth}{0pt}}
\caption{Free-run simulation on a held-out subject. Top: ground truth
versus NNARX-FROLS ($k=40$); the dominant oscillation is recovered.
Bottom: same ground truth versus the best LSTM, which collapses to a
near-zero trivial attractor.}
\label{fig:sim}
\end{figure}

\subsection{Kalman-assisted free-run simulation}\label{sec:kalman}

Section~\ref{sec:phase1} showed that the recurrent models collapse to
a trivial attractor in free run while the NNARX family saturates around
$\mathrm{VAF}\approx 27\%$. The natural next question is how much of
the residual error is due to error accumulation in the autoregressive
feedback, and how much to a genuinely insufficient model. To probe
this, we run a complementary experiment in which the closed-loop
feedback is no longer blind: at every step, a small fraction of the
true output is mixed back into the predicted output before it is fed
to the regressor (NNARX) or the recurrent cell. The companion notebook
\texttt{with\_simplified\_kalman.ipynb} contains the full
implementation; the bulk of the code is shared with the baseline
notebook, the only difference being the simulation routines.

\paragraph{Mechanism}
Let $\hat{y}[k]$ denote the model's prediction at step $k$ and
$y[k]$ the (would-be) sensor reading. Instead of feeding back
$\hat{y}[k-1]$ as in pure free-run, we feed back the
\emph{innovation-corrected} value
\begin{equation}
    \tilde{y}[k-1]
    = \hat{y}[k-1] + \alpha\,\big(y[k-1] - \hat{y}[k-1]\big),
    \label{eq:kalman}
\end{equation}
with a fixed scalar gain $\alpha\in[0,1]$. Equation~\eqref{eq:kalman}
is the steady-state form of a scalar Kalman filter with constant
process and measurement noise (also known as the
$\alpha$-$\beta$ filter~\cite{Brookner1998}); it can equivalently be
read as a linearly blended teacher-forcing scheme, with
$\alpha=0$ recovering pure free-run, $\alpha=1$ recovering full
one-step-ahead teacher forcing, and intermediate values producing a
\emph{soft} closed-loop simulation that requires only a fraction of
the true output to remain anchored. Throughout this experiment we use
$\alpha=0.15$, i.e.\ the simulator is allowed to peek at $15\%$ of the
sensor mismatch at every step.

\paragraph{Setup}
The model parameters are \emph{not} retrained: we take the same
checkpoints used in Section~\ref{sec:phase1} and only replace the
evaluation routine with the Kalman-assisted simulator. For the NNARX
family, $\tilde{y}$ replaces $y$ in the lagged regressor
\eqref{eq:nnarx}; for the recurrent models, $\tilde{y}[k-1]$ is
concatenated to $u[k]$ as input to the cell at step~$k$. The warm-up
window of $\max(n_u,n_y)$ samples (NNARX) and $10$ samples (recurrent)
is unchanged.

\paragraph{Results}
Table~\ref{tab:kalman} summarises the free-run accuracy of the same
models evaluated with and without the Kalman correction. The effect is
substantial. The best NNARX-FROLS variant jumps from
$\mathrm{VAF_{sim}}=27.2\%$ to $58.5\%$; the RNN, which essentially
collapsed in blind free-run ($2.7\%$), reaches $34.7\%$ with the
correction -- already above the NNARX-FROLS blind result of
Section~\ref{sec:phase1}.

\begin{table}[t]
\centering
\caption{Kalman-assisted free-run simulation, $\alpha=0.15$, LOSO split.
NRMSE/VAF are computed on the held-out subject. Rows marked
``LSTM/GRU not retested'' had their Kalman pass not executed in the
companion run; the corresponding cells in
\texttt{with\_simplified\_kalman.ipynb} are kept in place for
reproducibility.}
\label{tab:kalman}
\setlength{\tabcolsep}{4pt}
\begin{tabular}{lccc}
\toprule
\textbf{Model}
& \textbf{VAF\textsubscript{sim}}
& \textbf{VAF\textsubscript{sim}}
& \textbf{$\Delta$VAF} \\
& blind & $\alpha=0.15$ & \\
\midrule
NNARX-FROLS ($k=10$) & 25.9\% & 56.4\% & +30.5 \\
NNARX-FROLS ($k=20$) & 25.5\% & \textbf{58.5\%} & \textbf{+33.0} \\
NNARX-FROLS ($k=40$) & 27.2\% & 55.6\% & +28.4 \\
RNN                  &  2.7\% & 34.7\% & +32.0 \\
LSTM                 &  6.5\% & --     & (not retested) \\
GRU                  &  5.2\% & --     & (not retested) \\
\bottomrule
\end{tabular}
\end{table}

\paragraph{Interpretation}
The Kalman correction is not a fair head-to-head competitor of the
blind models -- it uses an additional input that pure simulation does
not have access to. We report it nonetheless because (i) it explains
why the blind recurrent models look so bad: they \emph{can} learn a
useful representation, but the closed-loop drift hides it; and
(ii) it is the exact configuration in which the present models would
be deployed if a (noisy) measurement of $y$ were available at
inference time, e.g.\ in a brain-computer-interface scenario in which
the EEG signal is being recorded online.

In other words, Table~\ref{tab:kalman} should be read as an
\emph{upper bound at $\alpha=0.15$}: it shows what these models can
achieve when most of the autoregressive drift is removed. The fact
that the NNARX-FROLS variant comfortably beats the published Volterra
baseline of~\cite{Vlaar2018} in this regime ($58\%$ vs $46\%$) is a
hint that the gap reported in Section~\ref{sec:discussion} is mostly
attributable to free-run drift, not to model capacity. A full sweep of
$\alpha$ and the corresponding LSTM/GRU runs are left to future work.

%==============================================================
\section{Discussion}\label{sec:discussion}

\subsection{Why does NNARX outperform the recurrent models in free run?}

The headline finding -- a plain MLP over a fixed regressor beats every
recurrent cell in closed-loop simulation -- is, at first sight,
counter-intuitive. There are three concurrent reasons.

\emph{First}, the NNARX is trained on a $k$-step curriculum up to
$k=320$ samples and is therefore explicitly exposed to its own
predictions during training. The recurrent models, even though they
ingest the input sequentially, are trained on a one-step MSE that does
not penalise long-horizon drift; scheduled
sampling~\cite{Bengio2015} and noise injection in the hidden state are
known mitigations but were not used here, to keep the protocol
uniform across families.

\emph{Second}, the dataset is small: ten subjects times seven
realisations of $30~\mathrm{s}$ each, after a $\times 8$ downsampling,
amount to a few minutes of effective training data. The NNARX
benefits from a very strong inductive bias (it explicitly receives
both $u[k\!-\!i]$ and $y[k\!-\!j]$ as inputs); the recurrent models
must learn to construct an equivalent regressor inside their hidden
state.

\emph{Third}, the recurrent training pipeline is sensitive to
back-propagation-through-time truncation, gradient clipping and warm-up
length. We used standard, fixed values for all three; tuning them on a
per-subject basis is a natural next step.

\subsection{Gap with the literature}

Our best free-run VAF ($27\%$) is roughly $20$ percentage points below
the Volterra baseline of~\cite{Vlaar2018} ($46\%$) and substantially
worse than the subject-specific NARMAX of~\cite{Gu2021}. The main
explanation is that Vlaar and Gu both fit
\emph{subject-specific} models, while we use a LOSO protocol that
treats the test subject as completely unseen; the literature numbers
are therefore an upper bound that the present setting is not allowed
to match. Subject-specific fine-tuning of the NNARX-FROLS model is the
most natural way to close the residual gap and is left to future work.

On one-step prediction the gap is much smaller: our
$\mathrm{VAF_{OSA}}\approx 85\%$ is within $10$ percentage points of
the $94\%$ reported by the best NARMAX and ensemble methods, despite
the much harder LOSO protocol.

\subsection{Limitations}
The study has several limitations worth stating explicitly. The
hyper-parameter search is coarse: only a handful of widths and depths
are explored per family. The recurrent training does not use
scheduled sampling, hidden-state noise injection, or per-subject
fine-tuning. Finally, although we report wall-clock training time,
energy and memory footprint are not measured.

%==============================================================
\section{Conclusion}\label{sec:conclusion}

We have compared four families of neural models for the
identification of the cortical response to wrist-joint perturbations,
on the publicly available Vlaar~2018 benchmark, using a unified
preprocessing pipeline and a single configuration-driven notebook.
On one-step ahead prediction every family is competitive
($\mathrm{VAF}\geq 75\%$); in closed-loop free-run simulation only the
NNARX family is usable, with the NNARX-FROLS variant peaking at
$\mathrm{VAF}\approx 27\%$ on a LOSO split, while RNN, LSTM and GRU
collapse below $7\%$.

The result is a clean, perhaps cautionary, message: on a small
biomedical dataset, an MLP over an explicitly designed regressor,
trained with a $k$-step curriculum, beats every recurrent architecture
we tried. A complementary Kalman-assisted simulation experiment
(Section~\ref{sec:kalman}) shows that the gap between the blind
free-run regime and the Volterra/NARMAX literature is dominated by
autoregressive drift rather than by model capacity: with as little as
$15\%$ innovation feedback, the NNARX-FROLS variant climbs to
$\mathrm{VAF}\approx 58\%$ and the RNN to $\mathrm{VAF}\approx 35\%$.
Closing the residual gap in pure free-run, and rescuing the recurrent
models with scheduled sampling and stability-aware losses, are the
two natural directions for follow-up work.

%==============================================================
\section*{Reproducibility}

All experiments are produced by two notebooks at the root of the
project repository. The headline experiments of
Sections~\ref{sec:phase1}--\ref{sec:phase3} are produced by
\texttt{project-si-mlicv.ipynb}, which contains a single global
configuration dictionary (\texttt{CFG}) and writes its outputs to
\texttt{phase1\_metrics.json}. The Kalman-assisted experiment of
Section~\ref{sec:kalman} is produced by
\texttt{with\_simplified\_kalman.ipynb}, which shares the same
\texttt{CFG} structure and differs only in the simulation routines
(NNARX \texttt{simulate} and recurrent \texttt{simulate\_seq}). The
dataset is the public benchmark of~\cite{Vlaar2018,VlaarDataset};
no proprietary data are used.

%==============================================================
\begin{thebibliography}{99}

\bibitem{Vlaar2018}
M.~P. Vlaar, T.~Solis-Escalante, A.~C. Schouten, F.~C.~T. van der Helm,
``Modeling the Nonlinear Cortical Response in EEG Evoked by Wrist
Joint Manipulation,'' \textit{IEEE Trans. Neural Syst. Rehabil. Eng.},
vol.~26, no.~1, pp.~205--215, Jan.~2018,
doi:\,10.1109/TNSRE.2017.2751650.

\bibitem{VlaarDataset}
M.~P. Vlaar \textit{et al.}, ``Benchmark dataset: Cortical
Responses Evoked by Wrist Joint Manipulation,'' 4TU.\,Centre for
Research Data, doi:\,10.4121/uuid:176d8f78-d9fd-491e-90e7-9370e249b701.

\bibitem{Gu2018}
Y.~Gu, T.~Solis-Escalante, A.~C. Schouten, F.~C.~T. van der Helm,
``A Novel Approach for Modeling Neural Responses to Joint
Perturbations Using the NARMAX Method and a Hierarchical Neural
Network,'' \textit{Frontiers Comput. Neurosci.}, vol.~12, art.~96, 2018,
doi:\,10.3389/fncom.2018.00096.

\bibitem{Gu2021}
Y.~Gu, T.~Solis-Escalante, A.~C. Schouten, F.~C.~T. van der Helm,
``Nonlinear Modeling of Cortical Responses to Mechanical Wrist
Perturbations Using the NARMAX Method,'' \textit{IEEE Trans. Biomed.
Eng.}, vol.~68, no.~3, pp.~948--958, 2021,
doi:\,10.1109/TBME.2020.3013545.

\bibitem{Santos2023}
G.~Santos \textit{et al.}, ``Decoding Electroencephalography Signal
Response by Stacking Ensemble Learning and Adaptive Differential
Evolution,'' \textit{Sensors}, vol.~23, no.~16, art.~7049, 2023.

\bibitem{Ljung1999}
L.~Ljung, \textit{System Identification: Theory for the User}, 2nd ed.,
Prentice Hall, 1999.

\bibitem{Nelles2001}
O.~Nelles, \textit{Nonlinear System Identification}, Springer, 2001.

\bibitem{Sjoberg1995}
J.~Sj\"oberg \textit{et al.}, ``Nonlinear black-box modeling in
system identification: a unified overview,'' \textit{Automatica},
vol.~31, no.~12, pp.~1691--1724, 1995.

\bibitem{Norgaard2000}
M.~N{\o}rgaard, O.~Ravn, N.~K. Poulsen, L.~K. Hansen,
\textit{Neural Networks for Modelling and Control of Dynamic Systems},
Springer, 2000.

\bibitem{Hochreiter1997}
S.~Hochreiter, J.~Schmidhuber, ``Long Short-Term Memory,''
\textit{Neural Computation}, vol.~9, no.~8, pp.~1735--1780, 1997.

\bibitem{Cho2014}
K.~Cho \textit{et al.}, ``Learning Phrase Representations using RNN
Encoder-Decoder for Statistical Machine Translation,''
\textit{Proc.\ EMNLP}, 2014.

\bibitem{Elman1990}
J.~L. Elman, ``Finding Structure in Time,''
\textit{Cognitive Science}, vol.~14, no.~2, pp.~179--211, 1990.

\bibitem{Bengio2015}
S.~Bengio, O.~Vinyals, N.~Jaitly, N.~Shazeer,
``Scheduled Sampling for Sequence Prediction with Recurrent Neural
Networks,'' \textit{Proc.\ NeurIPS}, 2015.

\bibitem{Andersson2019}
C.~Andersson, A.~H. Ribeiro, K.~Tiels, N.~Wahlstr\"om,
T.~B. Sch\"on, ``Deep Convolutional Networks in System Identification,''
\textit{IEEE CDC}, 2019.

\bibitem{Forgione2021}
M.~Forgione, D.~Piga, ``dynoNet: A neural network architecture for
learning dynamical systems,'' \textit{Int. J. Adapt. Control Signal
Process.}, vol.~35, no.~4, pp.~612--626, 2021.

\bibitem{Brookner1998}
E.~Brookner, \textit{Tracking and Kalman Filtering Made Easy}, Wiley,
1998.

\end{thebibliography}

\end{document}
